\section*{Hoofdstuk \ref{ch:b2:quadratic} --- Kwadratische vormen}

\begin{solution}{exo:b2:quadratic:1}
$q_1 = (x + 2y)^2 - 3y^2$: rang $2$, signatuur $(1, 1)$ (een vorm
van hyperbooltype).

$q_2$: maak het kwadraat in $x$ af: $q_2 = (x + y)^2 + y^2 + 2yz +
3z^2 = (x+y)^2 + (y + z)^2 + 2z^2$: rang $3$, signatuur $(3, 0)$
--- positief definiet.
\end{solution}

\begin{solution}{exo:b2:quadratic:2}
\omterm{def:b2:reduction:eigen}{Eigenwaarden} $2$ (op $\operatorname{Vect}(1,1,1)$) en $-1$ (op het
vlak $x + y + z = 0$). Orthonormaliseer: $e_1 =
\frac{1}{\sqrt3}(1,1,1)$; in het vlak geeft Gram--Schmidt op
$(1,-1,0), (1,0,-1)$ dat $e_2 = \frac{1}{\sqrt2}(1,-1,0)$, $e_3 =
\frac{1}{\sqrt6}(1,1,-2)$. Dan is $P = (e_1\ e_2\ e_3)$ orthogonaal
met $P^{\mathsf T}AP = \operatorname{diag}(2,-1,-1)$.

De vorm $q = 2xy + 2yz + 2zx$ heeft matrix $A$: in de gedraaide
coördinaten is $q = 2X^2 - Y^2 - Z^2$ --- hoofdassen; signatuur
$(1,2)$, in overeenstemming met
\cref{ex:b2:quadratic:gaussexample} (dezelfde vorm!).
\end{solution}

\begin{solution}{exo:b2:quadratic:3}
$u^{**} = u$: $\langle u^{**}x, y\rangle = \langle x, u^*y\rangle =
\langle ux, y\rangle$ voor alle $y$. $(uv)^* = v^*u^*$: $\langle
uvx, y\rangle = \langle vx, u^*y\rangle = \langle x,
v^*u^*y\rangle$. Kern: $y \in \ker u^* \iff \langle x, u^*y
\rangle = 0\ \forall x \iff \langle u(x), y\rangle = 0\ \forall x
\iff y \perp \operatorname{im} u$. Rangen: $\dim\ker u^* = n -
\operatorname{rk} u$ (orthogonaal complement), dus
$\operatorname{rk} u^* = \operatorname{rk} u$ volgens de
rangstelling --- de euclidische gedaante van de rangstelling voor
de \omterm{def:b2:linalg:transpose}{getransponeerde}.
\end{solution}

\begin{solution}{exo:b2:quadratic:4}
Spectraal: $A = PDP^{\mathsf T}$, met $D$ diagonaal met elementen
$\lambda_i$ die voldoen aan $\lambda_i^3 = \lambda_i$: $\lambda_i
\in \{-1, 0, 1\}$. Dan is $A^2 = PD^2P^{\mathsf T}$ met $D^2$
diagonaal met elementen $0/1$: $A^2$ is \omterm{def:b2:quadratic:adjoint}{symmetrisch} en idempotent
($(A^2)^2 = A^4 = A\cdot A^3 = A^2$) --- \omterm{def:b2:quadratic:adjoint}{symmetrisch} idempotent
$=$ orthogonale projectie (het is de projectie op
$\ker(A^2 - I) = \ker(A-I)\oplus\ker(A+I)$ langs $\ker A$, en deze
staan volgens de spectraalstelling loodrecht op elkaar).

In het algemeen is $P(A) = P\!\left(\text{diag}\right)$: $P(A)$
heeft dezelfde \omterm{def:b2:reduction:eigen}{eigenvectoren} en \omterm{def:b2:reduction:eigen}{eigenwaarden} $P(\lambda_i)$ --- de
``spectrale afbeelding'' op het \omterm{def:b2:reduction:diag}{diagonaliseerbare} niveau.
\end{solution}

\begin{solution}{exo:b2:quadratic:5}
Gesloten: $O(n) = g^{-1}(\{I\})$ voor de \omterm{def:b2:metric:continuity}{continue} $g(P) =
P^{\mathsf T}P$ (veeltermelementen). Begrensd: elke kolom van $P
\in O(n)$ is een eenheidsvector, dus alle elementen liggen in
$\intcc{-1}{1}$. Gesloten en begrensd in $\mathcal{M}_n(\R) \simeq
\R^{n^2}$: \omterm{def:b2:metric:compact}{compact} (\cref{thm:b2:metric:compactprops}~(2)).

Niet \omterm{def:b2:metric:connected}{samenhangend}: $\det$ neemt op $O(n)$ de twee waarden $\pm1$
aan, en een \omterm{def:b2:metric:continuity}{continue} surjectie op $\{-1, 1\}$ splitst de ruimte
(het argument van \cref{ex:b2:metric:glnr}).
\end{solution}

\begin{solution}{exo:b2:quadratic:6}
\emph{Bestaan:} $A = PDP^{\mathsf T}$ met $D =
\operatorname{diag}(\lambda_i)$, $\lambda_i \geq 0$; zet $B =
P\sqrt D P^{\mathsf T}$ met $\sqrt D =
\operatorname{diag}(\sqrt{\lambda_i})$: \omterm{def:b2:quadratic:adjoint}{symmetrisch}, positief
semidefiniet, $B^2 = A$.

\emph{Eenduidigheid:} zij $B$ \omterm{def:b2:quadratic:adjoint}{symmetrisch} positief semidefiniet met
$B^2 = A$. Dan commuteert $B$ met $A$; bijgevolg bewaart $B$ elke
\omterm{def:b2:reduction:eigen}{eigenruimte} $E_\lambda(A)$ (voor $Ax = \lambda x$: $A(Bx) = BAx =
\lambda Bx$). Op $E_\lambda(A)$ is de beperking van $B$
\omterm{def:b2:quadratic:adjoint}{symmetrisch} positief semidefiniet met kwadraat
$\lambda\,\mathrm{id}$; haar \omterm{def:b2:reduction:eigen}{eigenwaarden} $\mu$ voldoen aan $\mu^2
= \lambda$, $\mu \geq 0$: $\mu = \sqrt\lambda$ --- de beperking is
dus \omterm{def:b2:reduction:diag}{diagonaliseerbaar} met de enige \omterm{def:b2:reduction:eigen}{eigenwaarde} $\sqrt\lambda$ en
\emph{is} bijgevolg $\sqrt\lambda\,\mathrm{id}$. Omdat $E =
\bigoplus E_\lambda(A)$, is $B$ vastgelegd: $B = \sqrt A$.
\end{solution}

\begin{solution}{exo:b2:quadratic:7}
In een orthonormale eigenbasis is $\norm{Ax}_2^2 = \sum \lambda_i^2
x_i^2 \leq (\max_i \lambda_i^2)\norm x_2^2$, met gelijkheid in de
bijbehorende \omterm{def:b2:reduction:eigen}{eigenvector}: $\vertiii A_2 = \max\abs{\lambda_i}$.
Voor de gegeven matrix: \omterm{def:b2:reduction:eigen}{eigenwaarden} $3, -1$ (de tweelingzus van
\cref{ex:b2:quadratic:spectralexample}): $\vertiii A_2 = 3$.
\end{solution}

\begin{solution}{exo:b2:quadratic:8}
($\Rightarrow$) Het blok $A_k$ van formaat $k \times k$ linksboven
is de matrix van de beperking van de (definiete) vorm tot het
opspansel van de eerste $k$ basisvectoren: positief definiet, dus
haar \omterm{def:b2:reduction:eigen}{eigenwaarden} zijn positief en $\Delta_k = \det A_k > 0$.

($\Leftarrow$) Inductie naar $n$; $n = 1$ is duidelijk. Neem aan
dat alle $\Delta_k > 0$. Met inductie is $A_{n-1}$ positief
definiet: de vorm $q$ beperkt tot $F = \operatorname{Vect}(e_1,
\dots, e_{n-1})$ is definiet. Diagonaliseer $q|_F$ (Gauss):
coördinaten $y_1, \dots, y_{n-1}$ met $q|_F = \sum y_i^2$. In de
volledige ruimte geeft het afmaken van het kwadraat in de laatste
veranderlijke
\[
q = \sum_{i=1}^{n-1} \bigl(y_i + c_i x_n\bigr)^2 + c\,x_n^2
\]
voor geschikte constanten (verzamel de kruistermen in de
kwadraten). De reductie toont signatuur $(n-1 + \epsilon, \cdot)$
met $\epsilon$ de tekenbijdrage van $c$; en de \omterm{def:b2:linalg:det}{determinant} behoudt
onder \omterm{def:b2:quadratic:def}{congruentie} het teken van het product van de
diagonaalcoëfficiënten ($\det(P^{\mathsf T}AP) = (\det P)^2\det A$):
$\Delta_n > 0$ dwingt $c > 0$ af. Bijgevolg is $q$ een som van $n$
kwadraten van onafhankelijke vormen: positief definiet.
\end{solution}

\begin{solution}{exo:b2:quadratic:9}
Zij $(e_1, \dots, e_n)$ een orthonormale eigenbasis voor $\lambda_1
\geq \dots \geq \lambda_n$.

\emph{$\lambda_2 \leq$ het min-max:} voor elk hypervlak $H$ voldoet
de tweedimensionale $V = \operatorname{Vect}(e_1, e_2)$ aan
$\dim(H \cap V) \geq 1$ (Grassmann): kies een eenheidsvector $x \in
H \cap V$, $x = ae_1 + be_2$, $a^2 + b^2 = 1$:
\[
\langle u(x), x\rangle = \lambda_1 a^2 + \lambda_2 b^2 \geq
\lambda_2 :
\]
het maximum van elk hypervlak is $\geq \lambda_2$.

\emph{$\geq$:} voor $H = e_1^{\perp}$ heeft elke eenheidsvector $x
= \sum_{i\geq2} x_ie_i \in H$ dat $\langle u(x), x\rangle =
\sum_{i \geq 2} \lambda_i x_i^2 \leq \lambda_2$, aangenomen in
$e_2$: het maximum van dit hypervlak is precies $\lambda_2$. Het
minimum over $H$ is dus $\lambda_2$.
\end{solution}

\begin{solution}{exo:b2:quadratic:10}
\emph{Algebraïsche weg:} $2q = \bigl(\sum x_i\bigr)^2 - \sum
x_i^2$. Op het hypervlak $H : \sum x_i = 0$ (dimensie $n - 1$) is
$q = -\frac12\sum x_i^2$ negatief definiet; op de rechte $\R(1,
\dots, 1)$ is $q(t, \dots, t) = \binom n2 t^2 > 0$. Een deelruimte
waarop $q$ positief definiet is, snijdt $H$ triviaal en heeft dus
dimensie $\leq 1$: volgens Sylvester
(\cref{thm:b2:quadratic:sylvester}) is $s = 1$, en $t \geq n-1$
wegens $H$; rang $\leq n$ dwingt signatuur $(1, n-1)$ af, rang $n$.

\emph{Spectrale weg:} de matrix is $\frac12(J - I)$; $J$ heeft
\omterm{def:b2:reduction:eigen}{eigenwaarden} $n$ (op $(1,\dots,1)$) en $0$ (op $H$), dus heeft
$\frac12(J-I)$ \omterm{def:b2:reduction:eigen}{eigenwaarden} $\frac{n-1}{2}$ (eenmaal) en
$-\frac12$ ($n-1$ maal): één positieve, $n-1$ negatieve ---
dezelfde signatuur, volgens
\cref{cor:b2:quadratic:principalaxes}.
\end{solution}

\begin{solution}{exo:b2:quadratic:11}
Schrijf $B = C^{\mathsf T}C$ (\cref{exo:b2:quadratic:6} met $C =
\sqrt B$). Dan geldt, met $c_1, \dots, c_n$ de kolommen van
$C^{\mathsf T}$:
\[
\operatorname{tr}(AB) = \operatorname{tr}(AC^{\mathsf T}C)
= \operatorname{tr}(CAC^{\mathsf T})
= \sum_{i=1}^n \langle A c_i, c_i\rangle .
\]
Volgens \cref{cor:b2:quadratic:principalaxes}~(2) ligt elke term
tussen $\lambda_{\min}(A)\norm{c_i}^2$ en
$\lambda_{\max}(A)\norm{c_i}^2$, en $\sum\norm{c_i}^2 =
\operatorname{tr}(C^{\mathsf T}C)^{\vphantom1} =
\operatorname{tr} B$: de dubbele ongelijkheid volgt. Is $A$
bovendien positief semidefiniet, dan is $\lambda_{\min}(A) \geq 0$:
$\operatorname{tr}(AB) \geq 0$.
\end{solution}

\begin{solution}{exo:b2:quadratic:12}
\begin{enumerate}
  \item $\varphi(M, N) = \operatorname{tr}(MN)$ is bilineair en
        \omterm{def:b2:quadratic:adjoint}{symmetrisch} ($\operatorname{tr}(MN) =
        \operatorname{tr}(NM)$), en $\varphi(M, M) = q(M)$:
        $q$ is de \omterm{def:b2:quadratic:def}{kwadratische vorm} van $\varphi$.
  \item Voor \omterm{def:b2:quadratic:adjoint}{symmetrische} $S$ en antisymmetrische $K$ is
        $\operatorname{tr}(SK) =
        \operatorname{tr}\bigl((SK)^{\mathsf T}\bigr) =
        \operatorname{tr}(K^{\mathsf T}S^{\mathsf T}) =
        -\operatorname{tr}(KS) = -\operatorname{tr}(SK)$, dus
        $\varphi(S, K) = 0$: de twee deelruimten zijn
        $\varphi$-orthogonaal. Element voor element is
        $\operatorname{tr}(M^2) = \sum_{i,j} m_{ij}m_{ji}$: voor
        \omterm{def:b2:quadratic:adjoint}{symmetrische} $M$ is dit $\sum m_{ij}^2 > 0$ ($M \neq 0$);
        voor antisymmetrische $M$ is het $-\sum m_{ij}^2 < 0$.
  \item $\mathcal M_n(\R) = S_n \oplus A_n$ met dimensies
        $\frac{n(n+1)}2$ en $\frac{n(n-1)}2$; een \omterm{thm:b2:quadratic:gauss}{Gauss-reductie}
        aangepast aan deze $\varphi$-orthogonale splitsing
        schrijft $q$ als $\frac{n(n+1)}2$ positieve en
        $\frac{n(n-1)}2$ negatieve kwadraten: signatuur
        $\bigl(\frac{n(n+1)}2, \frac{n(n-1)}2\bigr)$
        (Sylvester), rang $n^2$: de vorm is niet-ontaard.
\end{enumerate}
\end{solution}

\begin{solution}{pb:b2:quadratic:1}
\textbf{1.} $G$ is \omterm{def:b2:quadratic:adjoint}{symmetrisch} wegens de symmetrie van het
inproduct, en
\[
X^{\mathsf T}GX = \sum_{i,j}x_ix_j\langle v_i, v_j\rangle
= \Bigl\|\sum_i x_iv_i\Bigr\|^2 \geq 0 ,
\]
met gelijkheid dan en slechts dan als $\sum x_iv_i = 0$: $G$ is
definiet dan en slechts dan als de enige nulcombinatie de triviale
is, dat wil zeggen dan en slechts dan als de familie onafhankelijk
is.

\textbf{2.} Met $B = \sqrt A$ (\cref{exo:b2:quadratic:6}): $A
= B^2 = B^{\mathsf T}B$, de Gram-matrix van de kolommen van $B$;
neem $C = B$. En $X^{\mathsf T}AX = \norm{CX}^2$, dus $A$ is
definiet dan en slechts dan als $CX \neq 0$ voor $X \neq 0$, dat
wil zeggen dan en slechts dan als $C$ inverteerbaar is.

\textbf{3.} De beperking van de vorm tot
$\operatorname{Vect}(e_i, e_j)$ heeft matrix
$\begin{pmatrix} a_{ii} & a_{ij}\\ a_{ij} & a_{jj}
\end{pmatrix}$, nog altijd positief semidefiniet: haar \omterm{def:b2:linalg:det}{determinant}
(het product van haar niet-negatieve \omterm{def:b2:reduction:eigen}{eigenwaarden}) is $\geq 0$:
$a_{ij}^2 \leq a_{ii}a_{jj}$. Dit is Cauchy--Schwarz voor de
vectoren $v_i, v_j$ van een Gram-realisatie.

\textbf{4.} \emph{Bestaan:} schrijf $A$ als de Gram-matrix van een
onafhankelijke familie $(v_1, \dots, v_n)$ (vragen 1--2).
Gram--Schmidt levert een orthonormale $(e_1, \dots, e_n)$ met
\[
v_k = \sum_{i \leq k} t_{ik}\,e_i,
\qquad t_{kk} = \bigl\| v_k - \operatorname{proj}_{k-1}v_k
\bigr\| > 0 ,
\]
zodat $T = (t_{ik})$ bovendriehoekig is met positieve diagonaal, en
\[
a_{jk} = \langle v_j, v_k\rangle
= \sum_i t_{ij}t_{ik} = (T^{\mathsf T}T)_{jk} .
\]
\emph{Eenduidigheid:} geldt $T_1^{\mathsf T}T_1 = T_2^{\mathsf
T}T_2$, dan voldoet $R = T_1T_2^{-1}$ aan $R^{\mathsf T}R = I$: $R$
is orthogonaal, en bovendien bovendriehoekig met positieve
diagonaal (product van zulke matrices). Dan is $R^{-1} = R^{\mathsf
T}$ tegelijk boven- (inverse van een bovendriehoekige) en
benedendriehoekig (\omterm{def:b2:linalg:transpose}{getransponeerde} van een bovendriehoekige):
diagonaal; een orthogonale diagonaalmatrix heeft elementen $\pm1$,
en de positiviteit dwingt $R = I$ af: $T_1 = T_2$.

\textbf{5.} Voor $i, j \leq k$ betreft $(T^{\mathsf T}T)_{ij} =
\sum_m t_{mi}t_{mj}$ alleen $m \leq \min(i,j) \leq k$: het leidende
blok van formaat $k\times k$ van $A$ is $T_k^{\mathsf T}T_k$, met
$T_k$ het leidende blok van $T$. Bijgevolg is $\Delta_k = (\det
T_k)^2 = (t_{11}\cdots t_{kk})^2$. Nu stuit de \omterm{thm:b2:quadratic:gauss}{Gauss-reductie} van
een positief definiete vorm in de natuurlijke volgorde nooit op een
verdwijnende kwadraatcoëfficiënt (de spillen zijn de
diagonaalelementen van de achtereenvolgens herleide positief
definiete blokken): zij levert $q = \sum_k d_k\ell_k^2$ met
$\ell_k = x_k + (\text{termen in } x_{k+1}, \dots)$, dat wil zeggen
$A = L^{\mathsf T}DL$ met $L$ unipotent driehoekig; dan is $T =
\sqrt D\,L$ een Cholesky-factor, dus wegens de eenduidigheid is
$t_{kk}^2 = d_k$ en
\[
d_k = \frac{(t_{11}\cdots t_{kk})^2}
{(t_{11}\cdots t_{k-1,k-1})^2}
= \frac{\Delta_k}{\Delta_{k-1}} .
\]
Op \cref{ex:b2:quadratic:tworoads}: $\Delta_1, \Delta_2,
\Delta_3 = 2, 3, 4$ en de spillen waren $2, \frac32, \frac43$.

\textbf{6.} Elke $a_{ii} = e_i^{\mathsf T}Ae_i > 0$. Zij $D =
\operatorname{diag}(a_{ii}^{-1/2})$ en $B = DAD$: positief definiet
(\omterm{def:b2:quadratic:def}{congruentie}), met $b_{ii} = 1$, dus $\operatorname{tr} B = n$.
Haar \omterm{def:b2:reduction:eigen}{eigenwaarden} $\mu_i > 0$ voldoen, volgens de ongelijkheid
tussen rekenkundig en meetkundig gemiddelde, aan
\[
\det B = \prod_i\mu_i
\leq \Bigl(\frac{\sum\mu_i}{n}\Bigr)^{\!n} = 1 ,
\]
en $\det B = (\det D)^2\det A = \dfrac{\det A}{\prod
a_{ii}}$: $\det A \leq \prod a_{ii}$.

\textbf{7.} De ongelijkheid tussen rekenkundig en meetkundig
gemiddelde is een gelijkheid dan en slechts dan als alle $\mu_i$
gelijk zijn (aan $1$); een \omterm{def:b2:quadratic:adjoint}{symmetrische} matrix met als enige
\omterm{def:b2:reduction:eigen}{eigenwaarde} $1$ is $PIP^{\mathsf T} = I$. Dus gelijkheid dan en
slechts dan als $B = I$, dat wil zeggen dan en slechts dan als
$a_{ij} = 0$ voor $i \neq j$: $A$ diagonaal.

\textbf{8.} Is $M$ singulier, dan zijn beide leden $\geq 0 =
\abs{\det M}$. Anders is $A = M^{\mathsf T}M$ positief definiet
met $a_{ii} = \norm{c_i}^2$ en $\det A = (\det M)^2$: vraag 6 geeft
$(\det M)^2 \leq \prod\norm{c_i}^2$. Gelijkheid dan en slechts dan
als $A = M^{\mathsf T}M$ diagonaal is (vraag 7), dat wil zeggen dan
en slechts dan als de kolommen paarsgewijs orthogonaal zijn.

\textbf{9.} $\abs{\det M}$ is het volume van het parallellepipedum
opgespannen door de kolommen: dat volume is hoogstens het product
van de ribbelengten, met gelijkheid precies voor rechthoekige
balken. Geldt $\abs{m_{ij}} \leq 1$, dan is $\norm{c_i} \leq
\sqrt n$, dus $\abs{\det M} \leq n^{n/2}$. (Het bereiken ervan
dwingt orthogonale kolommen met elementen $\pm1$ af: een
hadamardmatrix.)

\textbf{10.} $X^{\mathsf T}A^{\mathsf T}AX = \norm{AX}^2 > 0$ voor
$X \neq 0$ ($A$ inverteerbaar): $A^{\mathsf T}A$ is positief
definiet. Haar vierkantswortel $S$ is positief definiet
(\omterm{def:b2:reduction:eigen}{eigenwaarden} $\sqrt{\lambda_i} > 0$), dus inverteerbaar, en $Q =
AS^{-1}$ voldoet aan
\[
Q^{\mathsf T}Q = S^{-1}A^{\mathsf T}AS^{-1}
= S^{-1}S^2S^{-1} = I :
\]
$A = QS$ met $Q$ orthogonaal en $S$ positief definiet.

\textbf{11.} Geldt $A = QS = Q'S'$, dan is $S'^{\,2} =
S'^{\mathsf T}Q'^{\mathsf T}Q'S' = A^{\mathsf T}A = S^2$; twee
positief semidefiniete matrices met hetzelfde kwadraat vallen samen
(\cref{exo:b2:quadratic:6}): $S' = S$, en vervolgens $Q' = AS^{-1}
= Q$.

\textbf{12.} $\det(A + \varepsilon I)$ is een veelterm in
$\varepsilon$ die niet nul is: zij heeft eindig veel nulpunten, dus
een zekere rij $\varepsilon_k \to 0$ vermijdt ze. Schrijf $A +
\varepsilon_kI = Q_kS_k$ (vraag 10). $O(n)$ is \omterm{def:b2:metric:compact}{compact}
(\cref{exo:b2:quadratic:5}): een deelrij geeft $Q_{\varphi(k)}
\to Q \in O(n)$. Dan is
\[
S_{\varphi(k)} = Q_{\varphi(k)}^{\mathsf T}
\bigl(A + \varepsilon_{\varphi(k)}I\bigr)
\longrightarrow Q^{\mathsf T}A =: S,
\]
\omterm{def:b2:quadratic:adjoint}{symmetrisch} positief semidefiniet als limiet van zulke matrices
(gesloten voorwaarden), en $A = QS$. Bovendien is $S^2 = S^{\mathsf
T}S = A^{\mathsf T}QQ^{\mathsf T}A = A^{\mathsf T}A$, dus $S =
\sqrt{A^{\mathsf T}A}$ wegens de eenduidigheid. Voor singuliere $A$
is $S$ singulier en is $Q$ niet uniek: zij kan willekeurig gewijzigd
worden op $(\operatorname{im} S)^{\perp}$ --- uiterste geval $A =
0$, waar elke orthogonale $Q$ voldoet.

\textbf{13.} Diagonaliseer $S = P\Sigma P^{\mathsf T}$
(spectraalstelling), $\Sigma = \operatorname{diag}(\sigma_i)$ met
$\sigma_i \geq 0$ de \omterm{def:b2:reduction:eigen}{eigenwaarden} van $S = \sqrt{A^{\mathsf T}A}$.
Dan is
\[
A = QS = (QP)\,\Sigma\,P^{\mathsf T} = U\Sigma V^{\mathsf T},
\qquad U = QP,\ V = P \in O(n) .
\]

\textbf{14.} $\norm{Ax}^2 = x^{\mathsf T}S^2x \leq
\sigma_{\max}^2\norm x^2$ met gelijkheid in een bovenste
\omterm{def:b2:reduction:eigen}{eigenvector} van $S$: $\vertiii A_2 = \sigma_{\max}$ --- voor een
\omterm{def:b2:quadratic:adjoint}{symmetrische} matrix $A$ heeft $S = \sqrt{A^2}$ de
\omterm{def:b2:reduction:eigen}{eigenwaarden} $\abs{\lambda_i}$, waarmee \cref{exo:b2:quadratic:7} wordt
teruggevonden. \omterm{def:b2:linalg:det}{Determinant}: $\abs{\det A} =
\abs{\det U}\det\Sigma\abs{\det V} = \sigma_1\cdots\sigma_n$.
Bol: schrijven we $x = Vy$ met $\norm y = 1$, dan heeft $Ax =
U\Sigma y$ coördinaten $z_i = \sigma_iy_i$ in het orthonormale
assenstelsel van de kolommen van $U$: het beeld is $\{\sum
z_i^2/\sigma_i^2 = 1\}$, een ellipsoïde met halve assen
$\sigma_i$.

\textbf{15.} Een rotatie naar de hoofdassen
(\cref{cor:b2:quadratic:principalaxes}) maakt van $f$ de uitdrukking
$\lambda_1X^2 + \lambda_2Y^2 + \beta_1X + \beta_2Y + c$, met
$\lambda_1\lambda_2 = \det A$. Is $\det A \neq 0$, absorbeer dan de
lineaire termen met de translatie $X \mapsto X -
\frac{\beta_1}{2\lambda_1}$ (en evenzo $Y$): $\lambda_1X'^2
+ \lambda_2Y'^2 = c'$. Voor $\det A > 0$ (gelijke tekens): een
ellips ($c'$ met het juiste teken), een punt of leeg. Voor $\det
A < 0$: een hyperbool ($c' \neq 0$) of twee snijdende rechten. Is
$\det A = 0$ met rang $1$ (zeg $\lambda_2 = 0 \neq \lambda_1$):
$\lambda_1X'^2 + \beta_2Y + c''$, een parabool wanneer $\beta_2
\neq 0$; anders twee evenwijdige rechten, één rechte of leeg.
Niet-ontaarde gedaanten: ellips, hyperbool, parabool, bestuurd door
het teken van $\det A$.

\textbf{16.} Het kwadratische deel $x^2 + 4xy + y^2$ heeft matrix
$\begin{pmatrix}1 & 2\\ 2 & 1\end{pmatrix}$, \omterm{def:b2:reduction:eigen}{eigenwaarden} $3$ en
$-1$ met orthonormale richtingen $\frac{1}{\sqrt2}(1,1)$,
$\frac{1}{\sqrt2}(1,-1)$. In de gedraaide coördinaten $u =
\frac{x+y}{\sqrt2}$, $v = \frac{x-y}{\sqrt2}$ is $x^2 + y^2 = u^2
+ v^2$ en $2xy = u^2 - v^2$, dus de vorm is $3u^2 - v^2$, en $2x
- 2y = 2\sqrt2\,v$. De vergelijking wordt
\[
3u^2 - v^2 + 2\sqrt2\,v = 4
\quad\Longleftrightarrow\quad
3u^2 - \bigl(v - \sqrt2\bigr)^2 = 2 :
\]
een hyperbool met middelpunt in $(u, v) = (0, \sqrt2)$, dat wil
zeggen $(x, y) = (1, -1)$, met assen langs het gedraaide
assenstelsel.

\textbf{17.} $f(x) = (x - x_0)^{\mathsf T}A(x - x_0) + f(x_0)$
zodra $Ax_0 = -\frac b2$, dat wil zeggen $x_0 = -\frac12A^{-1}b$:
de gradiënt van $f$ verdwijnt precies daar ($x_0$ is het
symmetriemiddelpunt). Met $M = \begin{pmatrix} I & x_0\\ 0 &
1\end{pmatrix}$:
\[
M^{\mathsf T}\widetilde QM
= \begin{pmatrix}
A & Ax_0 + \frac b2\\[2pt]
\bigl(Ax_0 + \frac b2\bigr)^{\mathsf T} &
x_0^{\mathsf T}Ax_0 + b^{\mathsf T}x_0 + c
\end{pmatrix}
= \begin{pmatrix} A & 0\\ 0 & f(x_0)\end{pmatrix},
\]
en $\det M = 1$: $\det\widetilde Q = f(x_0)\det A$. De gecentreerde
vergelijking luidt $q(X) = -f(x_0)$: voor $f(x_0) = 0$ ontaardt zij
tot $q(X) = 0$ (twee rechten door het middelpunt als de signatuur
$(1,1)$ is, het enkele punt $x_0$ als $q$ definiet is); voor
$f(x_0) \neq 0$ is de kegelsnede een echte ellips of hyperbool.

\textbf{18.} $A^{-1} = -\frac13\begin{pmatrix} 1 & -2\\ -2 &
1\end{pmatrix}$ en $\frac b2 = (1, -1)$: $x_0 =
-A^{-1}\frac b2 = (1, -1)$, zoals gevonden in vraag 16. $f(x_0) =
q(1,-1) + 2 + 2 - 4 = (1 - 4 + 1) + 0 = -2 \neq 0$, en
$\det\widetilde Q = f(x_0)\det A = (-2)(-3) = 6 \neq 0$:
niet-ontaard; $\det A = -3 < 0$: een hyperbool --- en inderdaad
stemt de gecentreerde vergelijking $3u^2 - (v - \sqrt2)^2 =
-f(x_0) = 2$ overeen met vraag 16.

\textbf{19.} De matrix is $(1-\lambda)I + \lambda J$:
\omterm{def:b2:reduction:eigen}{eigenwaarden} $1 + 2\lambda$ (richting $(1,1,1)$) en $1 -
\lambda$ (dubbel, op $x + y + z = 0$). Gevallen:
\begin{itemize}
  \item $-\frac12 < \lambda < 1$: alle \omterm{def:b2:reduction:eigen}{eigenwaarden} positief:
        een omwentelingsellipsoïde om $(1,1,1)$ (een bol voor
        $\lambda = 0$);
  \item $\lambda = 1$: $q = (x+y+z)^2$: de vergelijking geeft de
        twee evenwijdige vlakken $x + y + z = \pm1$;
  \item $\lambda = -\frac12$: \omterm{def:b2:reduction:eigen}{eigenwaarden} $0, \frac32,
        \frac32$: een cirkelvormige cilinder met as $(1,1,1)$;
  \item $\lambda > 1$: signatuur $(1, 2)$: een tweebladige
        hyperboloïde;
  \item $\lambda < -\frac12$: signatuur $(2, 1)$: een
        eenbladige hyperboloïde.
\end{itemize}

\textbf{20.} De polaire vorm van $q$ is een inproduct
$\langle\cdot,\cdot\rangle_q$ op $E$. Bij vaste $x$ is $y \mapsto
\varphi'(x, y)$ (polaire vorm van $q'$) lineair, en dus gelijk aan
$\langle z_x, y\rangle_q$ voor een unieke $z_x$; $u(x) := z_x$ is
lineair (eenduidigheid), en $\langle u(x), y\rangle_q =
\varphi'(x,y) = \varphi'(y,x) = \langle u(y), x\rangle_q$: $u$ is
\omterm{def:b2:quadratic:adjoint}{symmetrisch} in de euclidische ruimte $(E,
\langle\cdot,\cdot\rangle_q)$. De spectraalstelling
(\cref{thm:b2:quadratic:spectral}) geeft een $q$-orthonormale
eigenbasis $(\varepsilon_i)$ met $u(\varepsilon_i) =
\mu_i\varepsilon_i$: daarin is $q(x) = \sum x_i^2$ en $q'(x) =
\langle u(x), x\rangle_q = \sum\mu_ix_i^2$.

\textbf{21.} Zij $P$ de matrix van die basis: \omterm{def:b2:quadratic:def}{congruentie} geeft
$P^{\mathsf T}AP = I$ en $P^{\mathsf T}BP =
\operatorname{diag}(\mu_i)$. Dan is
\[
\det(B - \mu A) = \det(P^{-\mathsf T})
\det\bigl(\operatorname{diag}(\mu_i) - \mu I\bigr)
\det(P^{-1})
= (\det P)^{-2}\prod_i(\mu_i - \mu) :
\]
de $\mu_i$ zijn de nulpunten van de bundel $\det(B - \mu A)$.

\textbf{22.} $\det(B - \mu A) =
\det\begin{pmatrix} -2\mu & 1-\mu\\ 1-\mu & -\mu\end{pmatrix}
= 2\mu^2 - (1-\mu)^2 = \mu^2 + 2\mu - 1$: nulpunten $\mu_\pm = -1
\pm \sqrt2$. Oplossen van $(B - \mu_\pm A)v = 0$ geeft $v_\pm = (1
- \mu_\pm,\ 2\mu_\pm)$ (de identiteit $(1-\mu)^2 = 2\mu^2$ in de
nulpunten bevestigt dit langs de tweede rij). De $A$-normen komen
er netjes uit: $q_A(v_\pm) = 2(1 + \mu_\pm^2)$, en men gaat na dat
$\varphi_A(v_+, v_-) = 0$ met behulp van $\mu_+ + \mu_- = -2$ en
$\mu_+\mu_- = -1$. De basis $\Bigl(\frac{v_+}{\sqrt{2(1 +
\mu_+^2)}}, \frac{v_-}{\sqrt{2(1+\mu_-^2)}}\Bigr)$ is orthonormaal
voor $A$ en diagonaliseert $B$ met elementen $\mu_\pm$.

\textbf{23.} Zou een inverteerbare $P$ beide vormen
diagonaliseren, dan zou de berekening van vraag 21 geven dat
$\det(B - \mu A) = (\det P)^{-2}\prod(d_{2i} - \mu d_{1i})$, een
reële veelterm die in reële lineaire factoren uiteenvalt. Maar hier
is
\[
\det(B - \mu A) = \det\begin{pmatrix} -\mu & 1\\ 1 &
\mu\end{pmatrix} = -\mu^2 - 1 ,
\]
van graad $2$ zonder reëel nulpunt: tegenspraak. (De $A$ met
lorentzsignatuur laat ``$B$-rotaties'' toe die geen reële assen
hebben.)

\textbf{24.} $A^{-1}B = A^{-1/2}\bigl(A^{-1/2}BA^{-1/2}\bigr)
A^{1/2}$ met $A^{1/2} = \sqrt A$ positief definiet
(\cref{exo:b2:quadratic:6}): $A^{-1}B$ is gelijkvormig met de
\emph{\omterm{def:b2:quadratic:adjoint}{symmetrische}} $A^{-1/2}BA^{-1/2}$, en dus \omterm{def:b2:reduction:diag}{diagonaliseerbaar}
met reële \omterm{def:b2:reduction:eigen}{eigenwaarden}. En $\det(B - \mu A) = \det A\cdot
\det(A^{-1}B - \mu I)$: de nulpunten $\mu_i$ van de bundel uit
vraag 21 zijn precies de \omterm{def:b2:reduction:eigen}{eigenwaarden} van $A^{-1}B$.

\textbf{25.} (i) Elk deel leunde op de spectraalstelling: via de
vierkantswortel (Cholesky, polair), de grenzen op de \omterm{def:b2:reduction:eigen}{eigenwaarden}
(Hadamard), de hoofdassen (kegelsneden) en de aan $q$ aangepaste
gedaante (simultane reductie). (ii) Gram-realisaties, Hadamard en
de \omterm{pb:b2:quadratic:1}{polaire ontbinding} overleven in de semidefiniete wereld; de
eenduidigheid van Cholesky, de spilformule en de simultane reductie
hebben definietheid nodig (de vragen 12 en 23 tonen precies hoe zij
falen). (iii) Gekoppelde kleine trillingen: de kinetische energie
(positief definiet) en de potentiële energie zijn twee \omterm{def:b2:quadratic:def}{kwadratische
vormen}; de basis van de vragen 20--22 vormt de \emph{normale modi}
van het systeem, en de $\mu_i$ zijn de kwadraten van de
hoekfrequenties.
\end{solution}
